\section{Construction and Hierarchy Configuration}
\label{sec:construction}
\subsection{Batched Construction of Irregular Graphs}
The base and overlays share two computational tasks: constructing a local
vector graph for a variable-size member set, and selecting cross edges for
an irregular pair of member sets.  Launching a GPU operation for every tiny
group repeats preparation and launch costs.  Conversely, a few large blocks
can dominate local graph construction.  Our pipeline separates host
metadata, memory-bounded GPU batches, and final host adjacency so that each
cost remains visible.

\begin{figure}[t]
\centering
\begin{tikzpicture}[font=\footnotesize,>=Latex,
 b/.style={draw,rounded corners,align=center,text width=2.05cm,
 minimum height=0.95cm,inner sep=4pt}]
\node[b] (meta) at (0,1.45) {Host partitions\\and topology};
\node[b] (desc) at (2.65,1.45) {Group/pair\\descriptors};
\node[b] (gpu) at (5.30,1.45) {GPU candidate\\and top-$k$ batches};
\node[b] (out) at (5.30,-0.25) {Compact IDs\\to host adjacency};
\node[b] (valid) at (2.65,-0.25) {Serialize, reload,\\quality evaluation};
\draw[->] (meta)--(desc);
\draw[->] (desc)--(gpu);
\draw[->] (gpu)--node[right]{bounded transfer}(out);
\draw[->] (out)--(valid);
\end{tikzpicture}
\caption{Construction dataflow. The GPU accelerates declared numerical
stages; topology and output adjacency retain host ownership. Structural
validation and query-quality evaluation are distinct checks.}
\label{fig:gpu-construction}
\end{figure}

\paragraph{Local graphs.}
The hybrid hierarchy builder handles small blocks with exact top-$R$ graphs (or the configured
bounded-complete alternative), medium blocks with sampled CPU Vamana, and large blocks
with GPU neighbor-candidate refinement followed by diversity pruning and
degree repair.  The fixed thresholds are 2048 and 8192 points, with 256
sampled candidates for the middle route.  The CUDA candidate path uses a
FastGrnnd-style adaptation of NN-Descent~\cite{dong2011nndescent}; pruning
follows relative-neighborhood principles~\cite{toussaint1980rng}.
These implementation settings are shared by compared structures.  A separate
full-GPU profile sends all blocks above the small completion threshold to
CUDA.  Neither route consults query Recall.

\paragraph{Cross edges.}
Each permitted group or block transition induces nearest-target tasks.
Descriptors identify source and target member sets, dimensions, and degree
budgets.  A conventional GPU path computes SGEMM distance tiles and runs
top-$k$ separately; the fused path evaluates distances and selects target
identifiers without materializing the complete score matrix in host memory.
This uses the same distance/selection bottleneck exploited by GPU similarity
search~\cite{johnson2021gpu}, specialized to many irregular tasks.  Vectors
and norms remain resident when they fit the configured device budget;
otherwise bounded chunks are streamed.  Only identifiers are transferred
when downstream adjacency does not consume distances.

\begin{algorithm}[t]
\caption{Construction of one declared hierarchy plan}
\label{alg:gpu-build}
\begin{algorithmic}[1]
\Require vectors, label trie, plan $(T_\ell,\tau_\ell)$, degree and memory budgets
\For{each level $\ell$ independently}
 \State partition direct members and build the declared block topology
 \For{each block $b$}
  \State select its declared CPU/GPU local-graph route by member count
  \State generate candidates, prune, and apply the degree-repair policy
 \EndFor
 \State form cross-edge task descriptors from the block topology
 \For{each descriptor batch fitting device memory}
  \State run the declared CPU or GPU distance/top-$k$ backend
  \State append compact identifiers to host adjacency with level owners
 \EndFor
 \State record effective backend, fallbacks, time, and allocated memory
\EndFor
\State serialize, reload, and validate membership and edge ownership
\State evaluate resulting index quality separately at the declared Recall targets
\end{algorithmic}
\end{algorithm}

The base build uses the same task separation on exact groups.  In the
accelerated base profile, optimized host metadata, CUDA local graphs, fused
cross edges, and a configured CPU additional-edge completion path jointly
determine construction time.  Consequently a full-builder speedup cannot be
attributed to fusion alone.

\subsection{Costs and Quality Contract}
The abstract uncovered-mass partition costs $O(m|\mathcal T|)$ for $m$
overlays, plus emitted membership traversal.  The actual builder additionally
sorts labels and member identifiers and constructs bitmaps, ancestor caches,
and dense ownership maps.  Prefix topology has linear edge count; the current
generic overlay LNG builder tests source, target, and possible intermediate
sets, giving an $O(B_\ell^3 h)$ worst-case bound for $B_\ell$ roots of maximum
label length $h$.  This implementation bound is different from its quadratic
maximum output size.

For an exact cross-edge batch with tasks $(U,V)$, distance arithmetic is
\begin{equation}
 \Theta\!\left(d\sum_{(U,V)}|U|\,|V|\right).
\end{equation}
The number of selected identifiers is at most
$\sum_{(U,V)}|U|\min(C,|V|)$ for target budget $C$.  Fusing selection reduces
intermediate traffic and launch overhead; it does not remove the underlying
distance work.  Approximate local-graph construction depends on candidate
counts and refinement rounds, so we report actual work and time rather than
assign it an exact-search bound.

Loaded hierarchy storage includes direct memberships and vector adjacency,
plus $O(m(N+G))$ dense point/group ownership maps.  Shared vector coordinates
need not be duplicated per overlay.  Device residency or streaming must fit
the declared memory budget, including candidates and output buffers.  Query
workspace also includes per-level visited state, seed buffers, and queue backing storage.

Numerical backend execution, successful serialization, and edge-ownership
checks establish that an index was built as requested.  They do not establish
equal ANN quality.  In particular, approximate candidate generation or TF32
arithmetic can change adjacency.  A speed-quality claim therefore requires
re-querying each accelerated index, including the fastest hybrid build.
The construction table reports measured time and resources; the current
paper does not infer a quality-preserving speedup from those columns alone.

\subsection{A Query-Calibration-Free Instantiation}
\label{sec:drh}
Independent overlays introduce depth and scale choices.  DRH uses a simple
structural proxy $T+N/T$, whose continuous minimum is $\sqrt N$, to choose
the first threshold.  It then derives a growth ratio from local degree $R$
and cross degree $C$:
\begin{equation}
 T_1=2^{\lfloor\log_2\sqrt N+1/2\rfloor},\quad
 \rho=\max(2,\mathrm{round}(R/C)),\quad T_{\ell+1}=\rho T_\ell.
\end{equation}
It materializes scales while $N/T_\ell\ge C$.  Ignoring the implementation
depth cap and integer-overflow guard, the derived overlay count is
\begin{equation}
 m=\max\!\left(0,1+\left\lfloor\log_\rho\frac{N}{CT_1}\right\rfloor\right).
\end{equation}
An overlay uses LNG when $N/T_\ell>R$ and prefix topology otherwise.  The
coarsest graph is thus sparsified once the estimated block count fits within
one local neighborhood budget.  Here $N/T_\ell$ is a proxy, not the measured
block count; threshold partitioning leaves residuals and can emit oversized
blocks.  Neither the proxy objective nor the topology rule is a theorem about
query runtime.

\begin{algorithm}[t]
\caption{DRH structure and optional mass gate}
\label{alg:drh}
\begin{algorithmic}[1]
\Require vector count $N$, local degree $R$, cross degree $C$
\State $T\gets\mathrm{nearestPow2}(\sqrt N)$;
  $\rho\gets\max(2,\mathrm{round}(R/C))$
\While{$N/T\ge C$ and depth cap is not reached}
 \State append $(T,\mathrm{LNG})$ if $N/T>R$, otherwise $(T,\mathrm{Trie})$
 \State $T\gets\rho T$
\EndWhile
\State \Return overlay plan and next-scale gate $T_{\mathrm{gate}}=T$
\end{algorithmic}
\end{algorithm}

For Amazon, $N=602{,}453$, $R=64$, and $C=4$ give
the plan $1024{:}\mathrm{LNG}$ and $16384{:}\mathrm{Trie}$.
DRH-v1 enables the multilayer backend when an authorized physical level
$h\ge2$ exists (level 2 in the evaluated plans); DRH-v2 also requires
sufficient direct mass at the highest such level, using the
next unmaterialized threshold.  The held-out policy experiment uses a
conservative LNG base.  It does not automatically choose between LNG and Trie
bases, and must not be substituted for the matched-provider factorial
winners.  The policy reads no query frequencies, latency, or Recall, but its
per-query authorization and mass gate naturally depend on $Q$.

\subsection{Maintenance Scope}
Inserting a label set changes a canonical trie path and may change nearby
terminal-prefix relations.  This avoids maintaining all minimal-superset
relations across unrelated branches at the label-plane level.  It does not
make full-index updates free: label postings and ancestor caches change,
mass thresholds may trigger repartitioning, and local/cross-block vector
graphs require repairs.  A deletion can split a block or invalidate portals.
The evaluated implementation is a static builder; incremental publication,
concurrent query visibility, and update amortization remain unmeasured design
questions.  We claim no dynamic-index result from this structural observation.
